\section{The near-BPS corner of the generalized cusp}
\label{sec:corner}

The continuum short-time law of the flipped loop is $-\sum_c\Gamma(\pi-\alpha_c,\pi)\log(\ell\mu)$ with
$\alpha_c\approx|s|/3$, so its coupling dependence is that of the generalized cusp anomalous dimension
$\Gamma(\varphi,\theta)$ in the corner $\varphi=\pi-u$, $\theta=\pi-xu$, $u\to0$. Two exact inputs control it.

\subsection{Exact slope at the BPS line}

At every loop order $\Gamma$ is a polynomial in $\cos\theta$ vanishing on the BPS line $\cos\theta=\cos\varphi$:
$\Gamma=\sum_{k\ge1}(\cos\theta-\cos\varphi)^kR_k(\varphi)$. The CHMS near-BPS formula \cite{CHMS}, derived from
the latitude matrix model, determines $R_1$ exactly:
$R_1(\varphi)=-2\varphi B(\tilde\lambda)/((1-\varphi^2/\pi^2)\sin\varphi)$ with $\tilde\lambda=\lambda(1-\varphi^2/\pi^2)$.
In corner variables $\cos\theta-\cos\varphi\approx-(1-x^2)u^2/2$, which must be written as
$-2\sin\frac{\theta+\varphi}2\sin\frac{\theta-\varphi}2$ in any numerical evaluation, $\sin\varphi\approx u$ and
$\tilde\lambda\approx2\lambda u/\pi$, so
\begin{equation}
\Gamma_1=(1-x^2)\frac{\pi^2}2B\Big(\frac{2\lambda u}\pi\Big)(1+O(u)):
\end{equation}
the slope at the BPS line, $-\partial_x\Gamma|_{x=1}=\pi^2B(2\lambda u/\pi)$, is exact at all couplings and depends
on the single effective coupling $g=\lambda u$.

\subsection{Two loops and all orders}

The planar two-loop cusp of Drukker and Forini \cite{DrukkerForini}, transcribed and validated by three
independent limits (small angles, the BPS-line slope at every $\varphi$, the antiparallel logarithm
$\lambda^2/8\pi^3$), gives in the corner
\begin{equation}
\Gamma(\pi-u,\pi-xu)=(1-x^2)\Big[\frac g{16\pi}-\frac{g^2}{192\pi^2}\Big]+O(gu)+O(g^2u\log u).
\end{equation}
Through $O(g^2)$ this equals $(1-x^2)(\pi^2/2)B(2g/\pi)$. Hence $H:=\lim_{u\to0}\Gamma(\pi-u,\pi)/u=\lambda(1-1/N^2)/16\pi$
has no two-loop correction, at finite $N$ as well since the two-loop cusp is pure $C_FC_A$.

\begin{proposition}[$H$ to all orders in perturbation theory]
Write $\xi=(\cos\theta-\cos\varphi)/\sin\varphi$ and $\Gamma^{(L)}=\sum_k\xi^kc_k^{(L)}(\varphi)$. Suppose that at $L$
loops the antiparallel singularities $\Gamma^{(L)}(\pi-u,\theta_j)=O(u^{-1}\,\mathrm{polylog}\,u)$ hold at $L-1$
distinct fixed $\theta_j$ (they hold at one and two loops and in the ladder limit). Then $c_k^{(L)}(\pi-u)=O(u^{k-1}\mathrm{polylog}\,u)$
for $k\ge2$, and $H(\lambda,N)=\lambda(1-1/N^2)/16\pi$ order by order in perturbation theory.
\end{proposition}

The bound follows by inverting a Vandermonde matrix; only $o(u^{-3})$ is needed. The statement is order by order: it
should not be called exact nonperturbatively, and the two-loop form $(1-x^2)(\pi^2/2)B(2g/\pi)$ should not be
extrapolated beyond two loops.

\subsection{Strong coupling: the classical corner}

For the planar classical string with ansatz $z=r\zeta(\hat\varphi)$ and an $S^5$ angle, the opening $\Omega=\pi-\varphi$,
the internal angle $\theta$ and $\Gamma/\sqrt\lambda$ are one-dimensional integrals over $(\zeta_0,r_1)$; the antiparallel
limit gives $\Gamma\Omega\to-4\pi^2/\Gamma(\tfrac14)^4$, the small-angle limit $B=\sqrt\lambda/4\pi^2$, and the curve
$r_1=1/\zeta_0^2$ is the BPS line $\theta=\pi/\sqrt{1+\zeta_0^2}=\pi-\Omega$ with $\Gamma\equiv0$. In the corner the
string is not thin: its depth is $\zeta_0\sim\sqrt\Omega$ while the opening is $\Omega$. With $\zeta_0=\sqrt{w\Omega}$,
\begin{equation}
x=\frac3{\pi w}-\frac{\pi w}4,\qquad \Gamma=\sqrt{\lambda\Omega}\,F,\qquad F=\frac14\Big(\sqrt w-\frac4{\pi^2w^{3/2}}\Big),
\end{equation}
valid for $2\lambda\Omega/\pi\gg1$, $\Omega\ll1$. $F$ vanishes at the BPS point $w=2/\pi$;
$F(0)=1/(\sqrt{2\pi}\,3^{3/4})=0.17501$; and $F'(1)=-\tfrac14\sqrt{2/\pi}$, which equals $-\pi^2B(2g/\pi)$ at large $g$:
the classical corner reproduces the one exactly known coefficient. $F$ is not proportional to $1-x^2$ (the $k=1$
part alone would give $F(0)\approx0.0997$). The physical branch $F(|x|)$ has a classical kink at $\theta=\pi$, where
antipodal points of $S^5$ are joined by an $S^4$ of great circles; quantum effects should smooth it. A byproduct is
the antiparallel-lines law near $\theta=\pi$ at strong coupling, $V_\parallel\approx-\sqrt\lambda(\pi-\theta)^{3/2}/(3^{3/2}\sqrt\pi L)$,
non-analytic in $\cos\theta$; we did not find it in the literature.

\subsection{Consequences for the flipped loop}

Per cusp $\Gamma_c=\Phi(\lambda\alpha_c,x=0)$ with $\alpha_c\approx|s|/3$. For $\lambda|s|\ll1$ the continuum law
$\log W\approx-(\lambda(1-1/N^2)|s|/24\pi)\log(\ell\mu)$ is one-loop exact perturbatively; for $1/\lambda\ll|s|\ll1$
(planar) $\log W\approx-0.2021\sqrt{\lambda|s|}\log(\ell\mu)$. Since $s=a\sqrt t$, the Loewner time sets the effective
coupling of the corner, $g(t)=\lambda|a|\sqrt t/3$: in a strongly coupled theory the continuum short-time law changes
from $|s|\log$ to $\sqrt{|s|}\log$ as $t$ grows past $\sim9/(\lambda a)^2$. The crossover function $\Phi(g,0)$ is not
determined beyond $O(g^2)$; the three-loop generalized cusp \cite{CHMSthree} would fix the next term.

\subsection*{Ledger}
\begin{center}\small
\begin{tabular}{@{}p{0.55\textwidth}p{0.40\textwidth}@{}}
\toprule
Statement & Status\\
\midrule
Exact $k=1$ part and BPS-line slope & Proof given CHMS eq.~8 (literature)\\
Two-loop corner; $H$ without two-loop term & Proof from the DF formulas, validated by three limits\\
$H$ one-loop exact to all orders & Order-by-order argument under the multi-angle bound (assumption)\\
Classical corner $F(x)$, $F(0)$, $F'(1)$ & Derived (leading order in $\Omega$), numerically checked; $\lambda\Omega\gg1$\\
$V_\parallel\propto(\pi-\theta)^{3/2}$ & Derived; not located in the literature\\
Loewner time as effective coupling & Corollary\\
\bottomrule
\end{tabular}
\end{center}
